% Version 3.7: normalization of the baseline logical EXPORT value.
% Include within Published results, after its introduction and before extraction.
\paragraph{Logical EXPORT values and source denotations.}
\label{sec:logical-carrier}
The baseline EXPORT interface specifies a logical tuple
\[
 E=(c_E,\beta_E,b_E,I_E,W_E):
\]
chain identifier, branch version, expected global base, ordered DATA
interval and its evidence. The evidence contains the regional
receipts, ancestry, boundaries, historical authorities and native
context required by that interface. These are structured logical
values. No deployed binary encoder is fixed by this specification.
The signed objects retain their original bytes and hashes under
normalization; it does not reserialize signed content.

The logical structure provides field and list positions for DATA and
DRAIN bodies, signed proposals, certificates, ancestry, admission
witnesses, boundaries and source proofs. A complete result is the
DATA and direct or covering receipt with the required validation
objects specified in \cref{sec:delivery,sec:admission}. Its occurrence
in these old tuple fields is determined before normalization is
defined. It can occur inside ancestry rather than in the sender's
declared output list. A wrong outer value or an invalid unrelated
proof does not change which objects those fields contain.

A reference retains its baseline denotation: the same content and
authenticated source identity that the original validation interface
uses. It is an edge to an object, not the object's bytes. Public source
objects can be represented inline or by aliases in the authenticated
retained dictionary. That dictionary retains authenticated raw
Espresso objects, including evidence from rejected records, and
authenticated L1 objects and proofs obtained through their own source
interface. Checkpoint objects are limited to the objects established
by that checkpoint's specified authentication contract.

An alias may identify an object by its content commitment or existing
authenticated source identity. An exact occurrence position may help
retrieve it, but a redundant explicit position is not a new validity
requirement. If several authenticated occurrences contain the same
object, derive a canonical representative from those occurrences.
An L1 authority or proof need not have been embedded in an earlier
Espresso record. The separate requirement that an Espresso selection
reference precede the containing EXPORT block remains unchanged.

Private caches implement this public object lookup; they do not add
new public denotations. A cache miss for an already specified object is
a retrieval task, while an absent required object is an incomplete
proof. A local regional cache reference can satisfy a point-to-point
validation bundle under its existing receiver-resolution condition;
that fact alone does not make a private-only object publicly supplied
by an EXPORT. This distinction follows the existing public
complete-result visibility obligation, rather than a new decoder's
success or failure.

\paragraph{Structural normalization.}
Write $\tagg{Flat}(E,\mathcal H)$ for the logical graph representation
of the tuple and its baseline authenticated source context $\mathcal H$.
It preserves every DATA and evidence field, including nested ancestry
and boundary evidence, in a fixed field/list order. Each typed object
keeps its original content, dependency edges and logical occurrence
path. Derive edges from the actual object fields, not a sender's
auxiliary index. Shared objects retain their incoming edges and
occurrence paths. References retain the same authenticated denotations
just described. This graph definition is not an instruction to fetch
every referenced target before processing the record.

The actual inventory is incremental. Traverse the supplied inline
structure and already retained authenticated objects, visiting each
actual object once. Preserve an unresolved reference as a typed edge;
do not wait for its target merely to finish this traversal. Extract
the required projection consisting of DATA bodies, FINAL signatures,
hash-linked paths and the fixed incoming context used by
$\tagg{CheckFINAL}$. Resolve references to these required objects through
the existing public input and retained-history interface. An optional
recovery authority, regional L1-reference proof or native sidecar is
not on this projection, so its acquisition cannot precede or gate
discovery of neighboring kernels. Inline structured children are still
inventoried even when their enclosing proof is invalid.

This distinction uses the baseline object meanings. An L1 authority
record denotes its specified authority fields, not a hidden regional
DATA/FINAL carrier; a native proof sidecar is not the missing DATA or
FINAL signatures. The complete-result transport interface supplies
the result's actual DATA and commitment evidence, with its required
ancestry, either in the logical record or through its existing public
source denotations. A kernel referenced in retained Espresso evidence
is read there. It is not obtained by first validating an optional
$\tagg{AUTH}_{r,g}$ or native claim. If a concrete compression hides
required DATA/FINAL bytes behind a new resolver, its compatibility and
access obligations must be proved separately; the logical graph lemma
does not manufacture them. No unresolved edge or cycle supplies a
missing required object.

Do not filter this traversal by outer chain, version, base or interval
hints, or stop it because a sibling or enclosing proof is invalid.
Each candidate still checks its own required objects. Preserve
alternative proofs of the same DATA separately so that one failed
proof does not erase another complete proof. The fixed logical
occurrence path, including a derived authenticated source occurrence
where needed, orders proof choices. A concrete framed encoding may
implement these paths but is not assumed by this definition.

\begin{proposition}[Preservation of logical result evidence]
\label{prop:carrier-preservation}
For a baseline logical EXPORT $E$ and its authenticated source context
$\mathcal H$, structural normalization preserves every complete
regional-result subgraph represented by its DATA and evidence fields,
with references interpreted under their unchanged baseline source
denotations. It therefore preserves that result's public FINAL kernel.
Changing only outer hints or adding unrelated invalid structured
evidence cannot remove the subgraph. Normalization creates no payload,
receipt, authority or append permission.
\end{proposition}
\begin{proof}
Fix the finite validation witness of a complete result in the original
logical fields. Induct over its object structure as a statement about
graph denotation. A leaf keeps its original content. For a compound
object every child remains represented independently of outer hints
and sibling validity; the induction hypothesis preserves the children
and their actual dependency edges. A shared object keeps its content
and additional occurrence edges. A reference denotes the same
authenticated source object before and after normalization, so the
corresponding case applies to that object. Retaining its unresolved
edge is sufficient for this representational statement; no acquisition
deadline for the target is inferred.
Thus every object and edge in the old witness survives. Its direct
receipt, DATA bodies and exact covering path in particular survive,
which preserves the public kernel of \cref{sec:receipt-kernel}.
For an actual complete public input, the DATA/FINAL projection supplied
by the inherited delivery interface is therefore the same projection
inventoried from its inline and retained source objects. Its extraction
does not require acquisition of the optional witness targets excluded
from $\tagg{CheckFINAL}$.
All signed content, authorities and original provenance are unchanged.
No traversal rule adds an absent leaf or uses an outer hint as an
inclusion condition. This proves both preservation and non-fabrication.
\end{proof}

This is a representation result. The revised interpreter deliberately
changes acceptance of bad-hint records; the proposition does not claim
that its acceptance trace equals the baseline's all-or-nothing EXPORT
trace. Finiteness of the inventory also supplies no processing-time
bound by itself.

\paragraph{Delivery scope and concrete encodings.}
The inherited Publish Layer and Espresso interfaces make the complete
logical result and its record available to the covered online,
history-following readers under their stated admission and load limits.
The original tuple fields and source denotations specify what is
delivered; \cref{prop:carrier-preservation} proves that normalization
preserves it. This does not define ``carrying'' to mean success of the
new extractor. A result already supplied by that public delivery
interface uses its original record even when its outer hints are
wrong. No corrected submission or acceptance-retry deadline is inferred.

A concrete encoder and decoder must preserve these logical values and
source denotations, including independently represented result objects
beside invalid hints or auxiliary evidence. The specification supplies
no blanket compatibility theorem for an unspecified legacy codec,
private-only resolver or arbitrary damaged bytes. Missing payload or
required ancestry was already excluded from complete results. Arbitrary
ciphertext without defined EXPORT semantics is not a logical tuple;
a documented old encoding of a tuple, however, requires a genuine
compatibility proof and cannot be dismissed as opaque.

The relevant-namespace visibility contract likewise is not bounded
discovery of arbitrary occurrences in unrelated namespaces. If the
phrase ``already recorded'' is given a broader implementation meaning
that discharges delivery for an occurrence outside the specified
public reader interface, an additional discovery argument is needed.
No such broader claim is used here. A wrong outer chain hint in a
record actually supplied through the existing interface is different:
the inventory still exposes its objects, whose signed domain determines
whether this interpreter may use them. Binary compatibility, external
access, native preparation and request inclusion are not consequences
of the structural lemma.
